% ============================================================
% section6_interactivity.tex - Раздел 6: Интерактивная визуализация
% ============================================================
% !TeX program = xelatex
% Компиляция: xelatex --shell-escape main.tex

\section{Интерактивная визуализация семейства спектров}\label{sec:interactivity}

В предыдущих разделах мы заложили полный инженерный фундамент: в Разделе~\ref{sec:database} мы спроектировали локальную СУБД с составными ключами и написали пакетный загрузчик \mycode{load_dataset_to_store}. В Разделе~\ref{sec:processing} мы научились пакетно обрабатывать все пары сигнал-фон и рассчитывать спектры. Теперь вся база данных полностью заполнена и готова к визуализации.

Главная задача этого раздела --- создать быстродействующий интерактивный графический интерфейс с сеткой графиков (plot grid), объединяющий отображение временного импульса и его спектра с мгновенным переключением углов, номеров повторений и ширины Гауссова окна.

\subsection{Навигация по графикам с сеткой (Plot Grid) и слайдером}

Перейдем к созданию графического интерфейса. Мы разработаем полноценную интерактивную панель, использующую концепцию \textbf{сетки графиков} (plot grid). Окно программы будет разделено на две части:
\begin{enumerate}
    \item \textbf{Верхняя панель} --- временной ТГц-сигнал образца (сырой и заокненный) с наложенной огибающей Гауссова окна.
    \item \textbf{Нижняя панель} --- соответствующий ему спектр пропускания в рабочем частотном диапазоне.
\end{enumerate}

Для полноценного анализа мы добавим три интерактивных элемента управления:
\begin{itemize}
    \item Слайдер для плавной регулировки полуширины временного фильтра Гаусса $\sigma$ в диапазоне от 5 до 300 пс.
    \item Радиокнопки для выбора угла установки первого поляризатора (угла поворота образца).
    \item Радиокнопки для выбора номера повторного измерения (реплики).
\end{itemize}

При взаимодействии с любым из элементов оба графика должны мгновенно перерисовываться без мерцания.

\subsubsection{Принципы интерактивности: Разметка vs Физика}

Прежде чем приступать к написанию кода, важно понять фундаментальный принцип работы графических интерфейсов. Студенту часто кажется сложным обилие строк, связанных с настройкой координат кнопок и ползунков. Однако весь интерактивный код четко разделяется на три логических блока:
\begin{enumerate}
    \item \textbf{Визуальные декорации (Разметка):} Функции вида \mycode{plt.axes([left, bottom, width, height])} просто размечают экран. Они указывают Matplotlib, в каких абсолютных координатах относительно краев окна (в долях от 0 до 1) нарисовать слайдер или панель кнопок.
    \item \textbf{Событийный цикл (Event Loop) и обратные вызовы (Callbacks):} Метод \mycode{on_changed} связывает графический элемент с вашей функцией (например, \mycode{update}). Matplotlib непрерывно слушает действия пользователя, и при изменении положения ползунка автоматически запускает привязанную функцию.
    \item \textbf{Вычислительное ядро (Физика):} Внутри функции \mycode{update()} происходит вся физика ЦОС --- расчет Гауссова окна для новой ширины, удаление постоянной составляющей, БПФ и деление спектров.
\end{enumerate}

\subsubsection{Профессиональное обновление данных в Matplotlib}

\warning{Никогда не вызывайте \texttt{ax.clear()} внутри функций обратного вызова (callbacks) интерактивных виджетов. Очистка осей сбрасывает все подписи, сетку, масштабы и приводит к сильному мерцанию интерфейса. Профессиональный подход --- построить линии один раз с помощью метода \texttt{ax.plot()} при инициализации, сохранить ссылки на объекты \texttt{Line2D}, а затем быстро обновлять их координаты по оси Y с помощью метода \texttt{line.set\_ydata()}.}

Полный код интерактивного модуля доступен в \textbf{\textit{Приложении~\ref{app:plots}}}. Рассмотрим ключевые моменты настройки виджетов управления и функции обновления данных:

\begin{listing}[H]
	\caption{Интерактивное управление сеткой графиков (\mycode{plots.py})}\label{lst:plot_grid}
	\begin{minted}{python}
# 6. Панели управления (Виджеты) слева
ax_radio_angle = plt.axes([0.02, 0.5, 0.16, 0.35], facecolor='#f5f5f5')
radio_angle = RadioButtons(ax_radio_angle, [f"Угол {a}°" for a in angles1])

ax_radio_rep = plt.axes([0.02, 0.22, 0.16, 0.18], facecolor='#f5f5f5')
radio_rep = RadioButtons(ax_radio_rep, [f"Повтор {r}" for r in reps])

ax_slider_sigma = plt.axes([0.25, 0.05, 0.5, 0.03], facecolor='#f5f5f5')
slider_sigma = Slider(ax_slider_sigma, 'Ширина окна (пс)', 5.0, 300.0, valinit=current_sigma, valstep=1.0)

# 7. Функция обновления графиков
def update(val):
    nonlocal current_angle, current_rep, current_sigma
    current_sigma = slider_sigma.val
    
    sig_k, bg_k = get_keys(current_angle, current_rep)
    if sig_k is None or bg_k is None:
        return
        
    t_s, E_s = data_store[sig_k]
    t_b, E_b = data_store[bg_k]
    
    E_s_dc = E_s - np.mean(E_s)
    E_b_dc = E_b - np.mean(E_b)
    
    p_idx = np.argmax(np.abs(E_b_dc))
    t_p = t_b[p_idx]
    
    w = np.exp(-0.5 * ((t_s - t_p) / current_sigma) ** 2)
    E_s_win = E_s_dc * w
    
    w_b = np.exp(-0.5 * ((t_b - t_p) / current_sigma) ** 2)
    E_b_win = E_b_dc * w_b
    
    spec_s = np.abs(rfft(E_s_win))
    d_t = t_s[1] - t_s[0]
    fr = rfftfreq(len(E_s_win), d=d_t)
    spec_b = np.abs(rfft(E_b_win))
    spec_b_s = np.maximum(spec_b, 1e-10)
    tr = (spec_s / spec_b_s) ** 2
    
    # Обновляем графики временного импульса
    line_bg.set_xdata(t_b)
    line_bg.set_ydata(E_b_win)
    line_sig.set_xdata(t_s)
    line_sig.set_ydata(E_s_win)
    line_win.set_xdata(t_s)
    line_win.set_ydata(w)
    
    ax1.set_title(f'Временной ТГц-импульс (Угол {current_angle}°, Повторение {current_rep})')
    ax1.relim()
    ax1.autoscale_view(scalex=True, scaley=False)
    
    # Обновляем графики спектров пропускания
    mask_new = (fr >= config.F_MIN) & (fr <= config.F_MAX)
    line_trans.set_xdata(fr[mask_new])
    line_trans.set_ydata(tr[mask_new])
    
    ax2.set_title(f'Спектр пропускания (Угол {current_angle}°, Повторение {current_rep})')
    ax2.relim()
    ax2.autoscale_view(scalex=False, scaley=True)
    
    fig.canvas.draw_idle()
	\end{minted}
\end{listing}

\subsubsection{Разбор кода: Интерактивное управление и обновление осей}

Разберем структуру и ключевые функции обновления графиков в листинге:

\begin{syntaxdesc}
    \item[Создание переключателей \mycode{RadioButtons(ax_radio_angle, ...)}] Инициализирует виджет переключения углов и повторных измерений в выделенных осях интерфейса.
    \item[Создание слайдера \mycode{Slider(ax_slider_sigma, ...)}] Настраивает слайдер с диапазоном $[5.0, 300.0]$ пс для регулировки ширины окна.
    \item[Локальное обновление \mycode{line_sig.set_ydata(E_s_win)}] Быстро обновляет значения по оси ординат для заокненного сигнала образца, исключая мерцание.
    \item[Обновление огибающей \mycode{line_win.set_ydata(w)}] Корректирует форму Гауссова окна на графике в реальном времени.
    \item[Запрос перерисовки \mycode{fig.canvas.draw_idle()}] Сообщает бэкенду о необходимости плавной перерисовки холста при следующем удобном цикле.
\end{syntaxdesc}

\note{Методы \texttt{ax.relim()} и \texttt{ax.autoscale\_view()} заставляют Matplotlib автоматически пересчитать новые границы графиков при смене данных. Мы разрешили автомасштабирование по оси X (\texttt{scalex=True}), но заблокировали его по оси Y (\texttt{scaley=False}) для временной области, так как диапазон Y жестко зафиксирован нами вручную от -1.5 до 1.5.}

Внешний вид нашей интерактивной двухпанельной панели представлен на \figref{fig:6_1}.

\smartfigure{fig6.1.png}{Двухпанельный графический интерфейс (plot grid) с переключателем RadioButtons для интерактивной инспекции сигналов и спектров.}{fig:6_1}

\subsection{Визуализация семейства спектров}

Помимо поштучной инспекции углов, физику важно увидеть общую картину --- семейство всех полученных спектров на одной координатной оси. Это позволит проследить плавную угловую модуляцию ТГц-излучения.

Напишем вспомогательную функцию \mycode{plot_all_transmissions()}, которая извлекает из СУБД все усредненные спектры и строит их, используя плавную цветовую карту градиента:

\begin{listing}[H]
	\caption{Построение семейства кривых пропускания (\mycode{plots.py})}\label{lst:plot_all_trans}
	\begin{minted}{python}
def plot_all_transmissions():
    data_store = data_loader.load_dataset_to_store(config.DATA_DIR)
    data_store = spectrum.process_dataset(data_store)

    spectra_keys = sorted([k for k in data_store.keys() if k[0] == 'spec_avg'])

    plt.figure(figsize=(10, 6))

    # Создаем красивый градиент цветов от фиолетового к оранжевому
    colors = plt.cm.plasma(np.linspace(0, 0.85, len(spectra_keys)))

    for key, color in zip(spectra_keys, colors):
        freqs, transmission = data_store[key]
        mask = (freqs >= config.F_MIN) & (freqs <= config.F_MAX)
        plt.plot(
            freqs[mask],
            transmission[mask],
            label=f"Угол {key[1]}°",
            color=color,
            linewidth=1.5
        )

    plt.xlim(config.F_MIN, config.F_MAX)
    plt.ylim(-0.05, 1.1)
    plt.xlabel("Частота (ТГц)")
    plt.ylabel("Коэффициент пропускания по мощности")
    plt.title("Семейство спектров пропускания для различных углов")
    plt.grid(True, linestyle='--', alpha=0.5)
    plt.legend(loc='lower left', ncol=2)
    plt.tight_layout()
    plt.show()
	\end{minted}
\end{listing}

Использование конструкции \texttt{plt.cm.plasma(np.linspace(...))} позволяет автоматически распределить цвета линий по спектральной карте. Это делает графики эстетичными и легко читаемыми. Итоговый график семейства кривых пропускания показан на \figref{fig:6_2}.

\smartfigure{fig6.2.png}{Семейство спектров пропускания образца при углах поворота поляризатора от 0 до 90 градусов.}{fig:6_2}

\begin{minitaskbox}
	\textbf{Мини-задание: Автоматизированное исследование ширины Гауссова окна} \\
	Проведи исследование влияния ширины фильтра Гаусса $\sigma$ с помощью созданного интерактивного интерфейса: \\
	\begin{enumerate}
		\item Запустите скрипт визуализации. Выберите угол $50^\circ$ и повторение 1.
		\item Установите слайдер \textbf{Ширина окна (пс)} в минимальное значение ($5.0$ пс). Обратите внимание на нижнюю панель со спектром пропускания. Почему кривая пропускания стала гладкой, а резкие провалы исчезли? Объясните эффект сглаживания спектра при сильном временном ограничении сигнала.
		\item Переместите слайдер в положение $30.0$ пс. Посмотрите на проступившие тонкие спектральные резонансы (линии поглощения паров воды). 
		\item Переместите слайдер на максимум ($300.0$ пс). Зафиксируйте появление периодических осцилляций (биений) и сильных высокочастотных шумов. Почему в спектр начали проникать паразитные шумы при избыточной ширине окна?
	\end{enumerate}
\end{minitaskbox}
\tip{Для быстрого создания аналитических веб-приложений (dashboards) в современной практике Data Science часто используют декларативный фреймворк \textbf{Streamlit}. В отличие от Matplotlib, в Streamlit вам не нужно настраивать сетки координат и прописывать функции обратного вызова (callbacks) --- веб-виджеты создаются одной строкой (например, \mycode{sigma = st.slider(...)}) и сразу возвращают выбранные значения в ваш линейный код, автоматически обновляя веб-страницу при перетаскивании ползунка.}

\subsection{Итоги раздела}

В этом разделе мы совершили качественный инженерный скачок в визуализации данных:
\begin{enumerate}
    \item Освоили инструмент сетки графиков (plot grid) в Matplotlib и научились мгновенно обновлять её с помощью высокопроизводительного обновления объектов \mycode{Line2D}.
    \item Построили интерактивное окно с плавной регулировкой ширины Гауссова окна с помощью виджета \mycode{Slider} в диапазоне $[5.0, 300.0]$ пс.
    \item Реализовали переключатели углов первого поляризатора и номеров повторений с помощью виджетов \mycode{RadioButtons}.
    \item Написали функцию \mycode{plot_all_transmissions()}, строящую всё семейство спектров пропускания на одном рисунке с красивым цветовым градиентом.
\end{enumerate}

\note{\textbf{Как изменился проект в этом разделе:}\\
В модуле \mycode{plots.py} добавлена логика интерактивного микро-приложения с RadioButtons, Slider и функцией \mycode{update}.\\
Также добавлена функция \mycode{plot_all_transmissions()} для построения семейства спектров.}

Теперь наш программный каркас полностью готов к глубокому физическому анализу. В следующем разделе мы подключим математический аппарат Blanco и научимся сопоставлять теорию с экспериментальным пропусканием на фиксированной частоте!
